\documentclass[conference]{IEEEtran}
\IEEEoverridecommandlockouts

\usepackage{cite}
\usepackage{amsmath,amssymb,amsfonts}
\usepackage{enumitem}
\usepackage[hidelinks]{hyperref}
\usepackage{algorithmic}
\usepackage{graphicx}
\usepackage{subcaption}
\usepackage{textcomp}
\usepackage{xcolor}
\usepackage{placeins}
\usepackage[most]{tcolorbox}
\def\BibTeX{{\rm B\kern-.05em{\sc i\kern-.025em b}\kern-.08em
    T\kern-.1667em\lower.7ex\hbox{E}\kern-.125emX}}

\newtcolorbox{takeaway}{
    colback=orange!15,
    colframe=orange!60!black,
    boxrule=0.5pt,
    arc=1pt,
    left=1pt, right=1pt, top=1pt, bottom=1pt,
    breakable
}
\begin{document}

\title{Characterizing DAOS for Scientific Streaming}


%%%%%% Author list (camera-ready byline; kept for after double-blind review, currently overridden below)
\author{
\IEEEauthorblockN{Xinxin Mei, Mark Jones, Vardan Gyurjyan}
\IEEEauthorblockA{
\textit{Department of Data and Computational Science} \\
\textit{Jefferson Lab}\\
Newport News, Virginia, USA \\
\{xmei, maj, gurjyan\}@jlab.org
% \orcid{0000-0003-1046-5269}
}
\and
\IEEEauthorblockN{Rebecca George$^{*}$\thanks{$^{*}$R. George conducted this work while affiliated with both William \& Mary and Jefferson Lab. She is no longer with either institution.}, Jie Ren}
\IEEEauthorblockA{
\textit{Department of Computer Science} \\
\textit{William \& Mary}\\
Williamsburg, Virginia, USA \\
\{rcgeorge, jren03\}@wm.edu
}
}


% },

% Double-blind PDSW submission: this second \author call overrides the byline above.
% \author{\IEEEauthorblockN{Anonymous Authors}}

\maketitle

\begin{abstract}
Distributed Asynchronous Object Storage (DAOS) is a high-performance, RDMA-based parallel storage system that we aim to use as a burst buffer for scientific streaming pipelines. Such pipelines require high concurrent throughput and low tail latency under continuous, line-rate data production. Prior DAOS evaluations emphasize aggregate throughput and scalability, while we additionally examine tail latency, read/write access behavior, and small-object performance. We characterize DAOS across two deployments, including a 128-storage-node allocation on Aurora. We find that I/O request size strongly affects throughput: MiB-scale requests approach client NIC bandwidth with modest concurrency, while small requests remain CPU- and IOPS-bound. DAOS also scales to multi-TiB/s across client nodes. At fixed total I/O concurrency, the process/queue-depth configuration substantially changes P99.99 completion latency while maintaining comparable throughput. Sequential and random throughput largely converge at large request sizes, while reading many small objects incurs substantial overhead at scale. These results provide configuration guidance for latency-sensitive scientific streaming systems.
\end{abstract}

\begin{IEEEkeywords}
DAOS, distributed storage system, performance benchmarking
\end{IEEEkeywords}


\section{Introduction}

Distributed Asynchronous Object Storage (DAOS) is a high-performance, open-source parallel I/O stack built on Non-Volatile Memory Express (NVMe) and Remote Direct Memory Access (RDMA) technologies~\cite{daos-exascale-storage2016sc, daos-intro2020, daos-dram-intro2023, daos-github}. It replaces conventional POSIX file and directory abstractions with a versioned, key-value object model and provides direct object and filesystem interfaces for high-performance applications. Since its initial deployment in 2020, prior studies have demonstrated strong DAOS scalability and performance, showing that it can match or outperform traditional parallel file systems across a variety of platforms and workloads~\cite{daos-vs-lustre2022pdsw, daos-arm2023SCAsia, daos-scalability2023SCAsia, daos-perf-explore2024scw, DAOS-aurora-20-node-test2024scw, XianheSun2023evaluation}. DAOS has also been studied for reducing I/O contention and accelerating data access in weather-prediction workflows, scientific data formats, and machine-learning pipelines~\cite{hdf5-daos2022tpds, daos-adios2024ISC, hpc-weather2023ipdps, manubens2024reducing, daos-smartnic2025scw, devarajan2021dlio, mlperf-storage-github, brim2025dug-daos-aio}. This work is motivated by using DAOS as a burst buffer for scientific streaming systems, where data from many scientific instruments arrive directly at high-performance computing (HPC) facilities and must be absorbed continuously at line rate without intermediate staging~\cite{ejfat2022indis}.

Scientific streaming places a different set of I/O demands on storage than the checkpoint- and training-oriented workloads emphasized in much of the prior DAOS literature. First, many instrument streams produce data concurrently, requiring high throughput under intra-node I/O concurrency and as the number of client nodes increases. Second, continuous data production makes latency important in addition to aggregate throughput. An individual slow write can delay a producer and introduce backpressure into the streaming pipeline, so mean latency alone does not capture the relevant behavior. Third, a streaming workflow spans different access phases, including write-dominated data production, mixed read/write online analysis, and read-dominated post-processing. Finally, raw detector events can be relatively small, approximately 1~MiB each in our target use case. The resulting performance therefore depends not only on aggregate bandwidth, but also on I/O request size, concurrency, access order, and how small events are organized into storage objects.


Existing DAOS evaluations establish its high peak throughput and scalability, but these results do not fully characterize how DAOS behaves under the combination of throughput, latency, and data-access requirements above. We therefore use controlled \textit{fio} and IOR experiments to isolate the storage parameters that are most relevant to scientific streaming. Rather than evaluating an end-to-end streaming application, our goal is to identify how these parameters affect DAOS performance and to derive configuration guidance for future streaming deployments. We characterize two different DAOS deployments, including a large-scale deployment on Aurora~\cite{aurora-intro2025}, and address the following research questions. \textbf{RQ1:} How do I/O request size, concurrency, DAOS access method, and client-node scale affect DAOS throughput? (Sec~\ref{sec:rq1_throughput}) \textbf{RQ2:} At fixed total I/O concurrency, how does the process/queue-depth configuration affect tail latency? (Sec~\ref{sec:rq2_latency}) \textbf{RQ3:} How do read/write mix, access order, and object size and count affect DAOS data access performance? (Sec~\ref{sec:rq3_us}) %We additionally compare DAOS with NFS to provide a conventional network file-system reference for request-size sensitivity. 

Our evaluation yields the following main findings:


\begin{comment}
    
\begin{itemize}
\item I/O request size relative to the DAOS chunk size is the first-order performance factor. Chunk-aligned requests saturate the NIC bandwidth with modest concurrency, whereas smaller requests remain IOPS- and CPU-bound.

\item At fixed total concurrency, the process/queue-depth split leaves throughput unchanged but can inflate P99.99 latency to 26$\times$ the median. More processes with shallower queues minimize tail latency.

\item DAOS scales near-linearly to multi-TiB/s across client nodes for large objects, but per-object overhead causes many-small-object reads to plateau at $3$--$5\times$ lower, motivating event aggregation before ingestion.
\end{itemize}


The remainder of this paper is organized as follows. Section~\ref{sec:background} reviews the DAOS architecture and prior performance studies. Section~\ref{sec:methods} describes our experimental platforms and benchmark configurations. Section~\ref{sec:results} presents our measurement results and findings. Section~\ref{sec:conclusion} concludes the paper and discusses future work.

\end{comment}






\begin{itemize}
\item I/O request size is the primary throughput factor under intra-node concurrency. MiB-scale requests approach client NIC bandwidth with modest concurrency, while small requests remain CPU- and IOPS-bound. On Aurora, DAOS scales to multi-TiB/s across client nodes.
\item At fixed total I/O concurrency and comparable throughput, the process/queue-depth configuration substantially changes P99.99 completion latency. Deep queues increase completion tail latency, while very shallow queues increase worst-case submission latency.
\item Sequential and random throughput largely converge at large request sizes, while small requests remain sensitive to access order. At scale, reading the same data volume from many small objects achieves 3--5$\times$ lower throughput than reading from large objects.
\end{itemize}
\section{Background and Related Work}

\subsection{DAOS Architecture}
Lofstead et al. presented the initial DAOS design as the next-generation HPC I/O stack by replacing POSIX file and directory abstractions with a distributed, object-based model featuring containers, versioned objects, and epoch-based consistency mechanisms \cite{daos-exascale-storage2016sc}.
Liang et al. followed with a full system description of DAOS built on top of persistent memory, the Persistent Memory Development Kit (PMDK), NVMe SSDs, and the Storage Performance Development Kit (SPDK) \cite{daos-intro2020}. Hennecke et al. extended the architecture to use DRAM as the primary storage tier in place of the persistent memory, demonstrating that DAOS maintains competitive performance on commodity hardware \cite{daos-dram-intro2023}.




% Jackson and Manubens evaluated DAOS as a general-purpose HPC storage layer, exploring its object, array, and POSIX-compatible interfaces under workloads typical of scientific computing \cite{daos-hpc-storage2023cluster}.

% The DAOS software stack is actively developed as an open-source project \cite{daos-repo}, and its client-side portability has been demonstrated on ARM64 platforms \cite{daos-arm2023SCAsia}.

% DAOS applications
% Using DAOS to speed up weather workflows \cite{daos2024weather_flow}. HDF5-DAOS \cite{hdf5-daos2022tpds}. Adios-DAOS \cite{daos-adios2024ISC}. DAOS as a burst buffer for EJFAT data streaming \cite{mei2024sc-poster}.

% \subsection{Aurora introduction}
Aurora introduction \cite{aurora-intro2025}. Slingshot RPC enhancement \cite{aurora-daos-slingshot-rpc2025cug}.

\subsection{DAOS Performance Analysis}
They also included performance results for a 10-node IO500 benchmarking \cite{daos-intro2020}. 
{
\color{blue}
(Needs some comparison between our results to their results.)
}

IOR benchmarking with server and client sides scaling \cite{daos-scalability2023SCAsia}.

The 20-node Aurora DAOS testbed results \cite{DAOS-aurora-20-node-test2024SC-W}. 

DAOS performance exploring \cite{daos-perf-explore2024SC-W}.

DAOS vs Lustre \cite{daos-vs-lustre2022pdsw}.


\section{Evaluation Methods}


\subsection{Experimental Platforms}

\subsubsection{Aurora Production}

\subsubsection{ Testbed}

\section{Evaluation Results}

\subsection{Single-Node Performance}

\section{Conclusions and Future Work}
\label{sec:conclusion}

\textcolor{blue}{
This paper presented a benchmarking study of DAOS across two production-representative platforms, Aurora's exascale-class deployment and a smaller DRAM/InfiniBand testbed, targeting the concurrency and access patterns of scientific streaming workloads.
Our results show that I/O concurrency behaves differently depending on access granularity: block size, rather than queue depth or process count, dominates single-process throughput, and large-block accesses saturate the network with far fewer processes than small-block accesses require. At fixed total concurrency, the split between per-process queue depth and process count strongly shapes tail latency even when mean throughput is unaffected, showing that concurrency \emph{structure}, not just its magnitude, is a first-order tuning knob for latency-sensitive streaming pipelines. At cluster scale, DAOS scales near-linearly to multi-TiB/s aggregate bandwidth with only tens of client nodes, and write throughput exceeds read throughput at the process counts typical of streaming ingestion.
}

\textcolor{blue}{
We further find that the choice of POSIX access method materially affects delivered performance: bare \textit{dfuse} mounts substantially underperform DAOS's native \texttt{dfs} interface, while the interception library \texttt{libpil4dfs} can recover, and even exceed, native performance for parallel \texttt{ior} workloads. Our read/write mix experiments show that DAOS writes consistently outperform reads, that small-block access is sensitive to access order (sequential access favors reads, random access favors writes), and that read throughput exhibits substantially higher run-to-run variability than write-dominated workloads. Finally, our comparison against NFS shows that DAOS's RDMA-based transport delivers an order-of-magnitude throughput advantage at or above the 1~MiB I/O size while consuming less CPU, though this advantage narrows, and can invert, once the workload becomes small-block and single-target IOPS-bound.
}

\textcolor{blue}{
Together, these findings offer concrete guidance for deploying DAOS in scientific streaming systems such as E2SAR: favor native \texttt{dfs} or interception-library access over bare POSIX, batch small I/O requests toward the 1~MiB DAOS chunk size where possible, and size process/queue-depth concurrency according to whether throughput or tail latency is the binding constraint.
}

\textcolor{blue}{
In future work, we plan to extend this study in three directions. First, we will profile server-side engine threads and PMem/DRAM-versus-NVMe tiering decisions as concurrency scales, to attribute the observed bottlenecks to DAOS-specific overheads rather than fabric- or device-level limits. Second, we will broaden our workload coverage beyond synthetic \textit{fio}/IOR benchmarks to end-to-end scientific streaming pipelines such as E2SAR, quantifying the gap between micro-benchmark peaks and realized application performance. Third, we will compare DAOS against additional storage backends and access layers, including kernel-bypass alternatives such as SPDK and NVMe-oF, to better delineate where DAOS's userspace, asynchronous, object-based design provides a genuine advantage over well-tuned conventional parallel file systems.
}

\section*{Acknowledgment}
We thank the anonymous reviewers for their valuable feedback. This research used resources of the Argonne Leadership Computing Facility, which is a U.S. Department of Energy Office of Science User Facility operated under contract DE-AC02-06CH11357. 
This material is based upon work supported by the U.S. Department of Energy, Office of Science, Office of Nuclear Physics under Contract No. 89243126CSC000213. The Jefferson Lab effort was also supported by the Laboratory Directed Research and Development Program at Thomas Jefferson National Accelerator Facility for the U.S. Department of Energy.
The William \& Mary effort was supported by NSF OAC-2348350.



% \newpage
\bibliographystyle{IEEEtran}
\bibliography{references}

% \appendix
% {\color{orange}

\section*{Interesting Questions for a DAOS Benchmarking Paper}

\subsection*{Performance Characterization}

\begin{itemize}
    \item How does DAOS's userspace I/O stack compare to kernel-bypass alternatives (SPDK-direct, NVMe-oF) at various concurrency levels?
    \item Where exactly does the bottleneck shift as you scale --- PMem bandwidth, NVMe throughput, network fabric, or DAOS engine threads?
    \item What is the overhead of DAOS's epoch/transaction model vs.\ raw device performance?
\end{itemize}

\subsection*{Workload Sensitivity}

\begin{itemize}
    \item How does the 4K--1M I/O size sweep expose DAOS's tiering decisions between PMem and NVMe?
    \item Does the read/write asymmetry change character across workload types (random, sequential, mixed)?
    \item How does object-level access (DAOS native API) compare to POSIX via DFS at the same block size?
\end{itemize}

\subsection*{Scalability and Concurrency}

\begin{itemize}
    \item Does performance degrade gracefully with many clients, or are there cliff edges?
    \item How does server-side thread count interact with client count --- is it a known bottleneck?
\end{itemize}

\subsection*{Practical HPC Relevance}

\begin{itemize}
    \item Does the benchmark reflect realistic application I/O patterns (checkpoint, analysis, streaming)?
    \item What is the gap between peak micro-benchmark numbers and what a real MPI application sees?
\end{itemize}

\subsection*{Core Question}

The question worth pushing hardest: \textbf{where does DAOS's architectural novelty
(PMem-first, userspace, asynchronous) actually show up in the numbers, vs.\ where it is
indistinguishable from a well-tuned Lustre or GPFS setup?}
That is the paper's core claim to make or break.

}


\end{document}
